% Figure rendered 5.5in wide with 8pt labels; included at 0.85\linewidth (6 columns).
% Same held-out pairs as 00_inpaper_fig_compare (Figure 3). The variants are Table 1's:
% morph-smooth = mnist/checkpoints/tune_mnist_smooth_g0.1_s0.3.
\begin{figure}[t]
\centering
\includegraphics[width=0.85\linewidth]{mnist_petct/00_inpaper_fig_ablation.pdf}
\caption{Ablation of MMCLAST-cg on MNIST-PET/CT, on the held-out pairs of
Figure~\ref{fig:mnist-compare}; the top three rows translate CT to PET and the
bottom three PET to CT. Columns show the source image, the target, and the
variants of Table~\ref{tab:mnist}: \emph{base} (Ours w/o Morph-smooth),
\emph{latcyc}, \emph{morph} and \emph{morph-smooth} (Ours w/ Morph-smooth).
CT$\rightarrow$PET has to invent the acquisition noise of the target: \emph{base},
\emph{latcyc} and \emph{morph-smooth} return a smooth, noise-free PET --- the
conditional mean --- while \emph{morph} reproduces the grain of the real PET and is
scored lower for it (Table~\ref{tab:mnist}). The grain follows the weight of the
endpoint adversarial loss rather than the critic $D_{\text{mix}}$:
\emph{morph-smooth} has no $D_{\text{mix}}$, yet with that weight at $1$, as in
\emph{morph}, instead of the selected $0.1$ it reproduces the grain too
(Figures~\ref{fig:mnist-compare} and~\ref{fig:mnist-path-ours}) and scores
$0.8405$, close to \emph{morph}'s $0.8387$. PET$\rightarrow$CT only has to remove
that noise, and the four variants are nearly indistinguishable.}
\label{fig:mnist-ablation}
\end{figure}
